%% Informe técnico — Detección de Deepfakes con M2F2-Det
%% Autores: Josu Zabala Muxika, Javier Solana Crespo
%% Mondragon Unibertsitatea — Trustworthy AI (ML2)
%%
\documentclass[pdflatex,sn-mathphys-num]{sn-jnl}

\usepackage{graphicx}
\usepackage{multirow}
\usepackage{amsmath,amssymb,amsfonts}
\usepackage{amsthm}
\usepackage{mathrsfs}
\usepackage[title]{appendix}
\usepackage{xcolor}
\usepackage{textcomp}
\usepackage{manyfoot}
\usepackage{booktabs}
\usepackage{algorithm}
\usepackage{algorithmicx}
\usepackage{algpseudocode}
\usepackage{listings}
\usepackage{hyperref}
\usepackage[utf8]{inputenc}
\usepackage[T1]{fontenc}
%% Nota: babel[spanish] omitido por incompatibilidad con sn-jnl \keywords{\and}
\renewcommand{\tablename}{Tabla}
\renewcommand{\figurename}{Figura}

\raggedbottom

\begin{document}

\title[Detección Multimodal de Deepfakes con M2F2-Det]{Detección Multimodal de Deepfakes mediante Modelos Visión-Lenguaje: Aplicación y Fine-Tuning de M2F2-Det}

\author[1]{\fnm{Josu} \sur{Zabala Muxika}}\email{josu.zabala@alumni.mondragon.edu}
\author[1]{\fnm{Javier} \sur{Solana Crespo}}\email{javier.solana@alumni.mondragon.edu}

\affil*[1]{\orgdiv{Facultad de Ingeniería}, \orgname{Mondragon Unibertsitatea}, \orgaddress{\city{Arrasate-Mondragón}, \state{Gipuzkoa}, \country{España}}}

\abstract{La proliferación de imágenes faciales sintéticas generadas por inteligencia artificial plantea retos crecientes para la verificación de la autenticidad visual. En este trabajo se estudia el modelo M2F2-Det (\emph{Multi-Modal Interpretable Forged Face Detector}), presentado como \emph{oral} en CVPR 2025, que combina el encoder multimodal CLIP, una red convolucional DenseNet-121 y un modelo de lenguaje grande (LLaVA-7B) para proporcionar tanto una clasificación binaria (real/falso) como una explicación textual de la decisión. Se describe el pipeline de inferencia de tres etapas implementado en Google Colab y se analiza experimentalmente su rendimiento sobre imágenes manipuladas por GAN, obteniendo un AUC-ROC del 100\,\% y una accuracy del 99,9\,\%. Sin embargo, se demuestra que el modelo original falla al evaluar imágenes generadas por difusión (Stable Diffusion, Grok/Aurora), diagnosticando las causas mediante análisis de las señales internas. Para resolver esta limitación se propone un fine-tuning con el dataset \emph{stable-diffusion-face-dataset}, compuesto por 9\,000 rostros generados con Stable Diffusion 1.5, 2.1 y SDXL. Se detallan la arquitectura, la metodología de reentrenamiento y los resultados obtenidos, enmarcando el trabajo dentro del campo de la IA Confiable (\emph{Trustworthy AI}).}

\keywords{Deepfake Detection, CLIP, Vision-Language Models, Trustworthy AI, Diffusion Models, Fine-Tuning}

\maketitle

%% ═══════════════════════════════════════════════════════════════════════
\section{Introducción}\label{sec:intro}
%% ═══════════════════════════════════════════════════════════════════════

El término \emph{deepfake} designa contenido multimedia ---en particular imágenes y vídeos de rostros--- generado o manipulado mediante técnicas de aprendizaje profundo con el objetivo de suplantar la identidad de una persona o crear contenido engañoso. Desde su aparición ligada a las redes generativas adversarias (GANs)~\cite{goodfellow2014generative}, los deepfakes han evolucionado rápidamente con la llegada de los modelos de difusión~\cite{rombach2022latentdiffusion}, capaces de producir imágenes de una calidad fotorrealista sin precedentes.

La detección automática de deepfakes es un pilar fundamental de la \textbf{IA Confiable} (\emph{Trustworthy AI}), ya que la desinformación visual puede afectar a procesos democráticos, perjudicar la reputación de personas y erosionar la confianza pública en los medios digitales. La Unión Europea, a través del AI Act, clasifica los sistemas de generación de contenido sintético como de alto riesgo, lo que subraya la necesidad de mecanismos de detección robustos y explicables~\cite{euaiact2024}.

Los detectores tradicionales de deepfakes se basan en redes convolucionales (CNNs) que aprenden artefactos de baja frecuencia propios de cada técnica de generación. Sin embargo, estos métodos presentan dos limitaciones fundamentales: (1)~generalización limitada a técnicas de generación no vistas durante el entrenamiento, y (2)~falta de explicabilidad, ya que producen una etiqueta binaria sin justificación. Para abordar ambos problemas, Guo et~al.~\cite{guo2025m2f2det} proponen \textbf{M2F2-Det}, un detector multimodal interpretable que combina el modelo visión-lenguaje CLIP~\cite{radford2021clip}, una red convolucional especializada y un modelo de lenguaje grande (LLM) para ofrecer simultáneamente detección y explicación textual.

En este trabajo se presenta la implementación práctica de M2F2-Det en Google Colab, se analiza su rendimiento sobre deepfakes basados en GAN, se diagnostica su fallo ante imágenes generadas por difusión y se propone un fine-tuning con datos de Stable Diffusion para mejorar su generalización. El informe se estructura como sigue: la Sección~\ref{sec:background} revisa los fundamentos teóricos (CLIP, BLIP, modelos de difusión); la Sección~\ref{sec:architecture} describe la arquitectura de M2F2-Det; la Sección~\ref{sec:methodology} detalla la metodología experimental; la Sección~\ref{sec:results} presenta los resultados; la Sección~\ref{sec:discussion} discute las implicaciones; y la Sección~\ref{sec:conclusion} ofrece las conclusiones.


%% ═══════════════════════════════════════════════════════════════════════
\section{Fundamentos Teóricos}\label{sec:background}
%% ═══════════════════════════════════════════════════════════════════════

\subsection{CLIP: Contrastive Language-Image Pre-training}\label{subsec:clip}

CLIP (\emph{Contrastive Language-Image Pre-training})~\cite{radford2021clip} es un modelo multimodal desarrollado por OpenAI que aprende representaciones conjuntas de imágenes y texto mediante aprendizaje contrastivo. Su arquitectura consta de dos encoders independientes ---un Vision Transformer (ViT)~\cite{dosovitskiy2020vit} para imágenes y un Transformer~\cite{vaswani2017attention} para texto--- entrenados sobre 400 millones de pares imagen-texto recogidos de Internet.

El objetivo de entrenamiento es maximizar la similitud coseno entre las representaciones de pares imagen-texto correspondientes y minimizarla para los no correspondientes. Formalmente, dada una imagen $I$ y un texto $T$, se computan:
\begin{equation}
    \mathbf{v} = f_{\text{vision}}(I), \quad \mathbf{t} = f_{\text{text}}(T)
\end{equation}
y la función de pérdida contrastiva (InfoNCE) opera sobre un batch de $N$ pares:
\begin{equation}
    \mathcal{L}_{\text{CLIP}} = -\frac{1}{N}\sum_{i=1}^{N} \log \frac{\exp(\text{sim}(\mathbf{v}_i, \mathbf{t}_i)/\tau)}{\sum_{j=1}^{N}\exp(\text{sim}(\mathbf{v}_i, \mathbf{t}_j)/\tau)}
    \label{eq:clip_loss}
\end{equation}
donde $\text{sim}(\cdot,\cdot)$ denota la similitud coseno y $\tau$ es un parámetro de temperatura aprendido.

En el contexto de la detección de deepfakes, CLIP aporta representaciones semánticas de alto nivel que capturan la relación entre la apariencia visual de un rostro y descripciones textuales como ``a real face'' o ``a manipulated face''. Sin embargo, como se mostrará, estas representaciones semánticas pueden sesgar la predicción hacia ``real'' cuando la imagen generada es visualmente coherente, como ocurre con los modelos de difusión.

La variante utilizada en M2F2-Det es \textbf{ViT-L/14@336}, que procesa imágenes a $336\times336$ píxeles con un Vision Transformer de 24 capas y produce embeddings de dimensión 1024.

\subsection{BLIP: Bootstrapping Language-Image Pre-training}\label{subsec:blip}

BLIP (\emph{Bootstrapping Language-Image Pre-training})~\cite{li2022blip} es un modelo visión-lenguaje que mejora el paradigma de CLIP mediante una estrategia de \emph{bootstrapping}: el modelo genera sus propios textos descriptivos (\emph{captions}) para imágenes web y utiliza un filtro para eliminar pares ruidosos, creando así un dataset de entrenamiento más limpio de forma iterativa.

A diferencia de CLIP, que solo realiza emparejamiento contrastivo, BLIP incorpora un decodificador de lenguaje que permite generar texto condicionado a la imagen, y un módulo de \emph{Image-Text Matching} (ITM) que clasifica si un par imagen-texto es coherente. Esta capacidad generativa es relevante para la explicabilidad en detección de deepfakes, ya que modelos derivados como BLIP-2~\cite{li2023blip2} sirven de base para LLMs multimodales como LLaVA.

Aunque M2F2-Det utiliza directamente CLIP como encoder visual y textual, la filosofía de BLIP influye en la segunda parte del pipeline: la generación de explicaciones textuales mediante un LLM (LLaVA), que hereda la arquitectura de alineación visión-lenguaje popularizada por BLIP-2.

\subsection{Modelos de difusión}\label{subsec:diffusion}

Los modelos de difusión~\cite{rombach2022latentdiffusion, ho2020ddpm} generan imágenes mediante un proceso iterativo de eliminación de ruido. A diferencia de las GANs, que producen la imagen en un solo paso a través de un generador, los modelos de difusión parten de ruido gaussiano puro y lo refinan progresivamente hasta obtener una imagen coherente.

\textbf{Stable Diffusion}~\cite{rombach2022latentdiffusion} opera en un espacio latente comprimido, lo que reduce el coste computacional manteniendo la calidad. Sus variantes (SD 1.5 a resolución $512\times512$, SD 2.1 a $768\times768$ y SDXL 1.0 a $1024\times1024$) producen imágenes con artefactos distintos a los de las GANs: mientras que las GANs dejan patrones espectrales característicos en frecuencias altas, los modelos de difusión generan texturas más suaves y naturales, lo que dificulta su detección por métodos entrenados exclusivamente con datos GAN.

Esta diferencia fundamental en los artefactos de generación es la motivación principal del fine-tuning propuesto en este trabajo.

\subsection{Redes convolucionales para detección de deepfakes}\label{subsec:cnn}

Las redes convolucionales profundas, en particular DenseNet~\cite{huang2017densenet} y EfficientNet~\cite{tan2019efficientnet}, se utilizan ampliamente como \emph{backbones} para la detección de deepfakes. Su capacidad para extraer características de textura a múltiples escalas las hace especialmente adecuadas para identificar artefactos sutiles.

\textbf{DenseNet-121}~\cite{huang2017densenet} conecta cada capa con todas las siguientes mediante conexiones densas, lo que favorece la reutilización de características y la propagación eficiente del gradiente. En M2F2-Det, DenseNet-121 actúa como \emph{deepfake encoder} especializado, extrayendo características de textura y frecuencia que complementan las representaciones semánticas de CLIP.



%% ═══════════════════════════════════════════════════════════════════════
\section{Arquitectura de M2F2-Det}\label{sec:architecture}
%% ═══════════════════════════════════════════════════════════════════════

M2F2-Det~\cite{guo2025m2f2det} es un detector de rostros falsificados multimodal e interpretable que opera en tres etapas: detección binaria (Stage~1), alineación multimodal (Stage~2) y generación de explicación textual (Stage~3).

\subsection{Stage 1: Detector binario}\label{subsec:stage1}

El detector binario combina tres ramas de procesamiento:

\begin{enumerate}
    \item \textbf{Rama CLIP-Vision}: El encoder visual ViT-L/14@336 (congelado) procesa la imagen y produce un token CLS ($\mathbf{v}_{\text{cls}} \in \mathbb{R}^{1024}$) y parches espaciales ($\mathbf{V}_{\text{patches}} \in \mathbb{R}^{576 \times 1024}$). El CLS se proyecta mediante una capa lineal y se escala por un parámetro aprendido $\alpha_{\text{vis}}$ (inicializado en 0,5).

    \item \textbf{Rama CLIP-Text}: El encoder textual procesa prompts de detección de deepfakes y produce embeddings de texto. Se computan puntuaciones de similitud coseno entre los parches visuales y los embeddings de texto, formando un vector de puntuaciones que se escala por $\alpha_{\text{txt}}$ (inicializado en 4,0).

    \item \textbf{Rama DenseNet-121}: La red convolucional (entrenable) extrae características de textura a partir de la imagen a $224\times224$. Los mapas de características intermedios de los bloques densos 1--3 se combinan con las características CLIP correspondientes mediante un \textbf{Bridge Adapter}, un módulo basado en Transformer que fusiona información de ambas ramas.
\end{enumerate}

Las tres salidas se concatenan para formar el vector final de características:
\begin{equation}
    \mathbf{f} = [\alpha_{\text{vis}} \cdot \mathbf{v}_{\text{cls}} \;\|\; \alpha_{\text{txt}} \cdot \mathbf{c}_{\text{bridge}} \;\|\; \text{proj}(\text{AvgPool}(\mathbf{d}))]
    \label{eq:feature_concat}
\end{equation}
donde $\mathbf{c}_{\text{bridge}} \in \mathbb{R}^{128}$ es la salida del Bridge Adapter, $\mathbf{d}$ son las características de DenseNet-121 y $\text{proj}(\cdot)$ es una proyección lineal. El vector $\mathbf{f} \in \mathbb{R}^{2 \times 1024 + 128}$ se pasa por una capa de salida $\text{Linear}(2176, 2)$ que produce logits para las clases ``real'' y ``falso''.

\subsection{Stage 2: Alineación multimodal}

En la segunda etapa, los embeddings del detector se alinean con el espacio de representación de un LLM (LLaVA-v1.5-7B~\cite{liu2023llava}) mediante un projector multimodal. Esta etapa entrena el projector para que el LLM pueda interpretar las señales del detector.

\subsection{Stage 3: Generación de explicación}\label{subsec:stage3}

La tercera etapa utiliza el modelo completo con el LLM para generar explicaciones textuales en lenguaje natural. Dado un prompt como:

\begin{quote}
\emph{``Determine the authenticity and explain your reasoning in detail. What visual cues indicate whether this face is real or manipulated?''}
\end{quote}

el modelo genera una respuesta que identifica la imagen como real o falsa y explica los indicios visuales (textura de piel, simetría, bordes, iluminación) que sustentan su decisión.

En nuestros experimentos con una imagen manipulada, Stage~1 la clasificó correctamente como \textbf{FAKE} con un 78,8\,\% de confianza. Sin embargo, Stage~3 (que utiliza el LLM) la clasificó erróneamente como real, generando una explicación detallada pero incorrecta. Esta discrepancia ilustra el sesgo semántico de los LLMs, que tienden a considerar ``reales'' las imágenes visualmente coherentes, y subraya la importancia de combinar múltiples señales para una detección robusta.


%% ═══════════════════════════════════════════════════════════════════════
\section{Metodología Experimental}\label{sec:methodology}
%% ═══════════════════════════════════════════════════════════════════════

La metodología experimental se divide en tres fases: (1)~implementación del pipeline de inferencia, (2)~entrenamiento del detector sobre datos GAN y evaluación, y (3)~diagnóstico de fallos con imágenes de difusión y propuesta de fine-tuning. Todo el trabajo se realizó en notebooks de Google Colab conectados desde VS Code, utilizando GPUs Tesla T4 y A100.

\subsection{Notebook de inferencia (M2F2\_Det\_Colab)}\label{subsec:colab}

El primer notebook implementa el pipeline completo de inferencia de M2F2-Det en 19 celdas organizadas en 5 secciones:

\begin{enumerate}
    \item \textbf{Setup} (Celdas 1--4): Verificación de GPU disponible, clonación del repositorio \texttt{CHELSEA234/M2F2\_Det}\footnote{\url{https://github.com/CHELSEA234/M2F2_Det}}, aplicación de parches de compatibilidad (cuantización 4-bit, fallback de \texttt{flash\_attn}, corrección de \texttt{map\_location}) e instalación de dependencias.

    \item \textbf{Descarga de pesos} (Celdas 5--6): Descarga automatizada de tres conjuntos de pesos: (a)~CLIP vision tower ($\sim$1,7\,GB), (b)~detector Stage~1 ($\sim$400\,MB), y (c)~modelo completo con LLM desde HuggingFace ($\sim$14\,GB).

    \item \textbf{Carga de imagen} (Celda 7): Montaje de Google Drive y carga de la imagen de test desde la ruta \texttt{ML2/test\_fotos/}.

    \item \textbf{Stage 1 --- Detección binaria} (Celda 8): Construcción del modelo \texttt{M2F2Det}, carga de pesos, preprocesamiento de la imagen a $224\times224$ y clasificación binaria con probabilidades softmax.

    \item \textbf{Stage 3 --- Detección + Explicación} (Celdas 9--12): Carga del modelo completo con LLaVA-7B en cuantización 4-bit (\texttt{load\_4bit=True}), aplicación de 7 parches de compatibilidad para ejecución en T4, y generación de la explicación textual.
\end{enumerate}

La ejecución sobre la imagen de test ``Manipulated-model.jpg'' (un rostro con una mitad retocada digitalmente versus la otra mitad natural) produjo los siguientes resultados:
\begin{itemize}
    \item \textbf{Stage 1}: Predicción \textbf{FAKE} con confianza del 78,8\,\%, P(Real)~=~21,2\,\%, P(Fake)~=~78,8\,\%.
    \item \textbf{Stage 3 (Detección)}: ``The image is real.'' --- predicción errónea.
    \item \textbf{Stage 3 (Explicación)}: El LLM generó una explicación detallada indicando cuatro evidencias de que la cara es real (piel suave natural, cejas arqueadas, ojos ovalados, diferencia de textura entre ambos lados). A pesar de ser una explicación coherente, la conclusión fue incorrecta.
\end{itemize}

\subsection{Notebook de fine-tuning (M2F2\_Det\_Finetune\_AIGen)}\label{subsec:finetune}

El segundo notebook (32 celdas) implementa el entrenamiento completo del detector Stage~1 y su fine-tuning posterior. Se organiza en las siguientes secciones:

\subsubsection{Secciones 1--8: Entrenamiento v1 (datos GAN)}\label{subsubsec:trainv1}

\begin{itemize}
    \item \textbf{Datos}: Se utiliza el dataset \texttt{saakshigupta/gradcam-fake-real-faces-xception-chunked}\footnote{\url{https://huggingface.co/datasets/saakshigupta/gradcam-fake-real-faces-xception-chunked}} que contiene rostros reales de FFHQ y rostros falsos generados por StyleGAN. Se equilibran las clases y se dividen en conjuntos de entrenamiento (80\,\%), validación (10\,\%) y test (10\,\%).

    \item \textbf{Modelo}: Se construye \texttt{M2F2Det} con DenseNet-121 como deepfake encoder (\texttt{hidden\_size=1024}), encoders CLIP ViT-L/14@336 congelados, y capa de salida $\text{Linear}(2176, 2)$.

    \item \textbf{Entrenamiento}: 10 épocas con optimizador Adam (lr~=~$10^{-4}$), \texttt{CrossEntropyLoss}, batch size de 32 y mixed precision (\texttt{float16}). Los mejores pesos se seleccionan por AUC-ROC en validación.
\end{itemize}

\subsubsection{Secciones 9--12: Evaluación y visualización}\label{subsubsec:eval}

Tras el entrenamiento se evalúa el modelo sobre el conjunto de test y se generan:
\begin{itemize}
    \item Métricas de clasificación (Accuracy, AUC-ROC, F1, Precision, Recall).
    \item Matriz de confusión.
    \item Curva ROC.
    \item Curvas de entrenamiento (loss y accuracy por época).
    \item Predicciones individuales con imagen y probabilidades.
    \item Mapas de activación \textbf{Grad-CAM}~\cite{selvaraju2017gradcam} sobre la última capa de DenseNet-121 (\texttt{denseblock4}), para visualizar qué regiones del rostro contribuyen más a la predicción.
    \item Test con imágenes de Google Drive (fotos propias y generadas por IA).
\end{itemize}

\subsubsection{Sección 13: Diagnóstico y fine-tuning con datos de difusión}\label{subsubsec:diagnostic}

Al evaluar el modelo entrenado con datos GAN sobre imágenes generadas con herramientas de difusión (Grok/Aurora de xAI), se observó que el detector las clasificaba consistentemente como \textbf{reales} con alta confianza. Para diagnosticar este fallo se implementó una celda de análisis de señales internas que evalúa independientemente la contribución de cada rama:

\begin{itemize}
    \item \textbf{Rama CLIP}: P(Real) = 99,9\,\% --- el encoder semántico considera la imagen como un rostro auténtico (esperado, ya que la difusión produce imágenes semánticamente coherentes).
    \item \textbf{Rama DenseNet}: P(Real) = 97,8\,\% --- la red convolucional tampoco detecta artefactos, porque fue entrenada solo con artefactos de GAN.
    \item \textbf{Predicción final}: P(Real) = 100,0\,\%.
\end{itemize}

Para abordar esta limitación se preparó un dataset de fine-tuning basado en imágenes de difusión:

\begin{itemize}
    \item \textbf{Dataset}: \texttt{tobecwb/stable-diffusion-face-dataset}\footnote{\url{https://github.com/tobecwb/stable-diffusion-face-dataset}} --- 9\,000 rostros generados con Stable Diffusion 1.5 (3\,000 a $512\times512$), SD 2.1 (3\,000 a $768\times768$) y SDXL 1.0 (3\,000 a $1024\times1024$), bajo licencia CC-BY-4.0.

    \item \textbf{Datos reales}: Se reutilizan 4\,500 imágenes reales del dataset de entrenamiento v1 (FFHQ).

    \item \textbf{División}: 4\,500 imágenes por clase, divididas en train (3\,600), validación (450) y test (450) por clase. Total: 9\,000 imágenes.
\end{itemize}

Las celdas de reentrenamiento (Celda 31) y evaluación del modelo reentrenado (Celda 32) fueron implementadas pero \textbf{no ejecutadas} en el momento de redacción de este informe, quedando como trabajo pendiente.


%% ═══════════════════════════════════════════════════════════════════════
\section{Resultados}\label{sec:results}
%% ═══════════════════════════════════════════════════════════════════════

\subsection{Entrenamiento v1: Datos GAN}\label{subsec:results_v1}

El entrenamiento sobre datos GAN (StyleGAN + FFHQ) convergió rápidamente, alcanzando métricas casi perfectas. La Tabla~\ref{tab:training_v1} muestra la evolución a lo largo de las 10 épocas.

\begin{table}[h]
\caption{Evolución del entrenamiento v1 (datos GAN) por época. Se muestra la loss y accuracy en entrenamiento (Train) y validación (Val), así como el AUC-ROC en validación.}\label{tab:training_v1}
\begin{tabular*}{\textwidth}{@{\extracolsep\fill}cccccc}
\toprule
Época & Train Loss & Train Acc & Val Loss & Val Acc & Val AUC \\
\midrule
1  & 0,2260 & 89,9\,\% & 0,0464 & 98,8\,\% & 99,93\,\% \\
2  & 0,0605 & 97,5\,\% & 0,0298 & 99,0\,\% & 99,93\,\% \\
3  & 0,0527 & 97,8\,\% & 0,0223 & 99,2\,\% & 99,97\,\% \\
5  & 0,0414 & 98,4\,\% & 0,0183 & 99,6\,\% & 99,98\,\% \\
7  & 0,0337 & 98,6\,\% & 0,0172 & 99,4\,\% & 99,99\,\% \\
10 & 0,0188 & 99,2\,\% & 0,0120 & 99,7\,\% & 99,99\,\% \\
\botrule
\end{tabular*}
\end{table}

\subsection{Evaluación en test: Datos GAN}\label{subsec:results_test}

Los resultados sobre el conjunto de test (1\,000 imágenes: 500 reales, 500 falsas GAN) se resumen en la Tabla~\ref{tab:test_v1}.

\begin{table}[h]
\caption{Métricas de evaluación en el conjunto de test v1 (datos GAN).}\label{tab:test_v1}
\begin{tabular*}{\textwidth}{@{\extracolsep\fill}ccccc}
\toprule
Accuracy & AUC-ROC & F1-Score & Precision & Recall \\
\midrule
99,9\,\% & 100,0\,\% & 99,9\,\% & 100,0\,\% & 99,8\,\% \\
\botrule
\end{tabular*}
\end{table}

La matriz de confusión fue:
\begin{itemize}
    \item 500/500 imágenes reales clasificadas correctamente.
    \item 499/500 imágenes falsas clasificadas correctamente.
    \item 1 falso negativo: una imagen GAN con apariencia atípica (rostro con rasgos exagerados tipo ``zombie'') que el modelo clasificó como real.
\end{itemize}

La curva ROC obtuvo un AUC de 1,0000, indicando separación perfecta entre las distribuciones de probabilidad de ambas clases.

\subsection{Visualización con Grad-CAM}\label{subsec:gradcam}

Los mapas Grad-CAM sobre la capa \texttt{denseblock4} de DenseNet-121 mostraron que el modelo concentra su atención en las siguientes regiones faciales:
\begin{itemize}
    \item \textbf{Zona periocular}: contorno de ojos y cejas, donde las GANs frecuentemente producen asimetrías.
    \item \textbf{Bordes del rostro}: transición entre la cara y el fondo, susceptible a artefactos de blending.
    \item \textbf{Textura de la piel}: regiones amplias de la frente y mejillas, donde se manifiestan patrones de frecuencia anómalos.
\end{itemize}

Estas activaciones son coherentes con la literatura sobre detección de deepfakes basados en GAN y confirman que el modelo aprendió características discriminativas significativas.

\subsection{Fallo en imágenes de difusión}\label{subsec:diffusion_failure}

Al evaluar el modelo con imágenes de prueba propias y generadas por herramientas de IA basadas en difusión, los resultados fueron significativamente distintos. La Tabla~\ref{tab:individual_tests} muestra las predicciones del modelo sobre imágenes individuales cargadas desde Google Drive.

\begin{table}[h]
\caption{Predicciones individuales del modelo v1 (entrenado con GAN) sobre imágenes de test heterogéneas.}\label{tab:individual_tests}
\begin{tabular*}{\textwidth}{@{\extracolsep\fill}lccc}
\toprule
Imagen & Tipo real & Predicción & Confianza \\
\midrule
Imagen generada por IA nº1 & AI-Gen & AI-Gen & 93,1\,\% \\
Imagen generada por IA nº2 & AI-Gen & AI-Gen & 100,0\,\% \\
Foto real nº1 & Real & Real & 100,0\,\% \\
Foto real nº2 & Real & Real & 53,3\,\% \\
\botrule
\end{tabular*}
\end{table}

Sin embargo, cuando se evaluaron imágenes generadas específicamente por modelos de difusión avanzados (Grok/Aurora de xAI), el modelo las clasificó como reales con confianza del 100\,\%. El diagnóstico de señales internas reveló que tanto la rama CLIP (P(Real) = 99,9\,\%) como la rama DenseNet (P(Real) = 97,8\,\%) conducían a esta clasificación errónea, confirmando la hipótesis de que los artefactos de difusión son cualitativamente diferentes de los de GAN.

Los parámetros $\alpha$ aprendidos durante el entrenamiento ($\alpha_{\text{vis}} = 0,6244$, $\alpha_{\text{txt}} = 4,4632$) indican que el modelo otorga mucho más peso a la rama textual (Bridge Adapter) que a la visual directa, lo que amplifica el sesgo semántico de CLIP.


%% ═══════════════════════════════════════════════════════════════════════
\section{Discusión}\label{sec:discussion}
%% ═══════════════════════════════════════════════════════════════════════

\subsection{Rendimiento excepcional en dominio GAN}

Los resultados del entrenamiento v1 demuestran que la arquitectura de M2F2-Det es altamente efectiva para detectar deepfakes basados en GAN. La combinación de características semánticas (CLIP) con características de textura (DenseNet-121), fusionadas mediante el Bridge Adapter, logra una separación prácticamente perfecta entre imágenes reales y falsas GAN. La convergencia rápida (AUC $>$ 99,9\,\% en la primera época) sugiere que los artefactos de GAN son fácilmente discriminables por esta arquitectura.

\subsection{Gap de dominio GAN vs. difusión}

El fallo ante imágenes de difusión revela un problema fundamental: \textbf{el detector aprende artefactos específicos de la técnica de generación, no características universales de artificialidad}. Las GANs producen patrones espectrales característicos (artefactos de checkerboard, inconsistencias en altas frecuencias) que DenseNet-121 aprende a detectar. Los modelos de difusión, en cambio, operan mediante refinamiento iterativo que produce texturas más naturales y artefactos diferentes (suavizado excesivo, inconsistencias de iluminación global).

El análisis de señales internas confirma que ambas ramas fallan:
\begin{itemize}
    \item \textbf{CLIP} falla porque su entrenamiento contrastivo sobre datos naturales le enseña representaciones semánticas, no forenses. Una imagen de difusión fotorrealista es semánticamente indistinguible de una foto real.
    \item \textbf{DenseNet-121} falla porque sus filtros convolucionales se especializaron en artefactos GAN durante el entrenamiento y no reconocen los patrones de difusión.
\end{itemize}

\subsection{Propuesta de fine-tuning con datos de difusión}

La solución propuesta ---reentrenamiento parcial con el dataset de Stable Diffusion--- aborda directamente la raíz del problema exponiendo a DenseNet-121 a artefactos de difusión. El dataset elegido incluye tres variantes de Stable Diffusion (1.5, 2.1, SDXL), lo que introduce diversidad en las resoluciones y arquitecturas de generación.

Este enfoque tiene limitaciones conocidas:
\begin{itemize}
    \item \textbf{Generalización}: El modelo reentrenado podría especializarse en Stable Diffusion y fallar ante otros modelos de difusión (DALL-E, Midjourney, Flux).
    \item \textbf{Catastrophic forgetting}: El reentrenamiento podría degradar el rendimiento sobre datos GAN si no se incluyen datos GAN en el nuevo dataset.
    \item \textbf{Tamaño del dataset}: 4\,500 imágenes por clase es relativamente pequeño para fine-tuning profundo, aunque suficiente para un ajuste del encoder convolucional.
\end{itemize}

\subsection{Explicabilidad: fortaleza y debilidad}

La capacidad de Stage~3 para generar explicaciones textuales es una fortaleza diferencial de M2F2-Det frente al paradigma de la IA Confiable. Sin embargo, nuestros experimentos revelaron que el LLM (LLaVA-7B) puede generar explicaciones plausibles pero incorrectas. La imagen manipulada fue explicada como real con argumentos coherentes sobre la textura de piel y la simetría facial. Esto plantea un riesgo de \emph{falsa confianza}: un usuario que reciba una explicación detallada y convincente podría confiar en una predicción errónea. Este hallazgo subraya la necesidad de combinar las salidas del detector binario (Stage~1) con las del LLM (Stage~3) y nunca confiar exclusivamente en la explicación textual.


%% ═══════════════════════════════════════════════════════════════════════
\section{Conclusiones y Trabajo Futuro}\label{sec:conclusion}
%% ═══════════════════════════════════════════════════════════════════════

En este trabajo se ha implementado y evaluado experimentalmente el modelo M2F2-Det para la detección multimodal e interpretable de deepfakes. Las principales contribuciones y hallazgos son:

\begin{enumerate}
    \item Se ha desplegado exitosamente el pipeline completo de M2F2-Det (Stage~1 a Stage~3) en Google Colab, incluyendo la resolución de múltiples problemas de compatibilidad (cuantización 4-bit, flash attention, device mapping).

    \item El modelo entrenado con datos GAN alcanza un rendimiento excepcional (AUC-ROC = 100\,\%, Accuracy = 99,9\,\%), confirmando la eficacia de la fusión CLIP + DenseNet-121 + Bridge Adapter.

    \item Se ha identificado y diagnosticado el \textbf{gap de dominio GAN-difusión}, demostrando que tanto la rama semántica (CLIP) como la convolucional (DenseNet) fallan ante imágenes de difusión.

    \item Se ha propuesto e implementado un pipeline de fine-tuning con 9\,000 imágenes de Stable Diffusion (1.5/2.1/SDXL), cuya ejecución y evaluación constituye el trabajo inmediato pendiente.

    \item Se ha observado que las explicaciones del LLM (Stage~3) pueden ser convincentes pero incorrectas, lo que plantea retos para la IA Confiable y sugiere la necesidad de mecanismos de calibración o verificación cruzada.
\end{enumerate}

\textbf{Trabajo futuro:}
\begin{itemize}
    \item Ejecutar y evaluar el fine-tuning v2 con datos de difusión.
    \item Evaluar la generalización a otros generadores (DALL-E 3, Midjourney, Flux).
    \item Investigar técnicas de \emph{continual learning} para evitar el olvido catastrófico al incorporar nuevos tipos de deepfakes.
    \item Calibrar la explicación del LLM con la confianza del detector binario para mejorar la fiabilidad de Stage~3.
    \item Fine-tuning del LLM (Stage~3) con datos de difusión para que las explicaciones reflejen los artefactos específicos de esta técnica.
\end{itemize}

%% ═══════════════════════════════════════════════════════════════════════

\backmatter

\bmhead{Agradecimientos}

Los autores agradecen a Mondragon Unibertsitatea y al programa de Trustworthy AI el apoyo proporcionado. Se reconoce el uso de Google Colab para las ejecuciones experimentales y los recursos de código abierto del proyecto M2F2-Det.

\section*{Declaraciones}

\begin{itemize}
    \item \textbf{Financiación}: Este trabajo se realizó en el marco de la asignatura de Machine Learning 2 de Mondragon Unibertsitatea.
    \item \textbf{Conflictos de interés}: Los autores declaran no tener conflictos de interés.
    \item \textbf{Disponibilidad de código}: El código de M2F2-Det está disponible en \url{https://github.com/CHELSEA234/M2F2_Det}. El dataset de Stable Diffusion en \url{https://github.com/tobecwb/stable-diffusion-face-dataset}.
    \item \textbf{Contribución de los autores}: Ambos autores contribuyeron igualmente a la implementación, experimentación y redacción del informe.
\end{itemize}

\begin{thebibliography}{99}

\bibitem{guo2025m2f2det}
Guo, X., Zhang, Y., Liu, J., Lyu, S.:
Rethinking vision-language model in face forensics: Multi-modal interpretable forged face detector.
In: Proceedings of the IEEE/CVF Conference on Computer Vision and Pattern Recognition (CVPR), pp.~1--12 (2025).
\url{https://arxiv.org/abs/2503.20188}

\bibitem{radford2021clip}
Radford, A., Kim, J.W., Hallacy, C., et~al.:
Learning transferable visual models from natural language supervision.
In: International Conference on Machine Learning (ICML), pp.~8748--8763. PMLR (2021)

\bibitem{li2022blip}
Li, J., Li, D., Xiong, C., Hoi, S.:
BLIP: Bootstrapping language-image pre-training for unified vision-language understanding and generation.
In: International Conference on Machine Learning (ICML), pp.~12888--12900. PMLR (2022)

\bibitem{li2023blip2}
Li, J., Li, D., Savarese, S., Hoi, S.:
BLIP-2: Bootstrapping language-image pre-training with frozen image encoders and large language models.
In: International Conference on Machine Learning (ICML), pp.~19730--19742. PMLR (2023)

\bibitem{liu2023llava}
Liu, H., Li, C., Wu, Q., Lee, Y.J.:
Visual instruction tuning.
In: Advances in Neural Information Processing Systems (NeurIPS), vol.~36 (2023)

\bibitem{goodfellow2014generative}
Goodfellow, I., Pouget-Abadie, J., Mirza, M., et~al.:
Generative adversarial nets.
In: Advances in Neural Information Processing Systems (NeurIPS), vol.~27 (2014)

\bibitem{rombach2022latentdiffusion}
Rombach, R., Blattmann, A., Lorenz, D., Esser, P., Ommer, B.:
High-resolution image synthesis with latent diffusion models.
In: Proceedings of the IEEE/CVF Conference on Computer Vision and Pattern Recognition (CVPR), pp.~10684--10695 (2022)

\bibitem{ho2020ddpm}
Ho, J., Jain, A., Abbeel, P.:
Denoising diffusion probabilistic models.
In: Advances in Neural Information Processing Systems (NeurIPS), vol.~33, pp.~6840--6851 (2020)

\bibitem{huang2017densenet}
Huang, G., Liu, Z., Van~Der~Maaten, L., Weinberger, K.Q.:
Densely connected convolutional networks.
In: Proceedings of the IEEE/CVF Conference on Computer Vision and Pattern Recognition (CVPR), pp.~4700--4708 (2017)

\bibitem{tan2019efficientnet}
Tan, M., Le, Q.:
EfficientNet: Rethinking model scaling for convolutional neural networks.
In: International Conference on Machine Learning (ICML), pp.~6105--6114. PMLR (2019)

\bibitem{dosovitskiy2020vit}
Dosovitskiy, A., Beyer, L., Kolesnikov, A., et~al.:
An image is worth 16x16 words: Transformers for image recognition at scale.
In: International Conference on Learning Representations (ICLR) (2021)

\bibitem{vaswani2017attention}
Vaswani, A., Shazeer, N., Parmar, N., et~al.:
Attention is all you need.
In: Advances in Neural Information Processing Systems (NeurIPS), vol.~30 (2017)

\bibitem{selvaraju2017gradcam}
Selvaraju, R.R., Cogswell, M., Das, A., et~al.:
Grad-CAM: Visual explanations from deep networks via gradient-based localization.
In: Proceedings of the IEEE International Conference on Computer Vision (ICCV), pp.~618--626 (2017)

\bibitem{euaiact2024}
European Parliament and Council of the European Union:
Regulation (EU) 2024/1689 laying down harmonised rules on artificial intelligence (Artificial Intelligence Act).
Official Journal of the European Union, L series (2024)

\bibitem{zhang2024ddvqa}
Zhang, Y., Guo, X., Liu, J., Lyu, S.:
Common sense reasoning for deepfake detection.
In: European Conference on Computer Vision (ECCV), pp.~1--17. Springer (2024)

\end{thebibliography}

\end{document}
